\documentclass[14pt,a4paper]{extarticle}
% =============================================================
% preamble.tex — оформление отчётов по СТП БГУИР 01-2024
% =============================================================

% Шрифты: под pdflatex — Times через tempora + T2A,
% под xelatex/tectonic — системный Times New Roman (запасной вариант STIX Two Text)
\usepackage{iftex}
\ifPDFTeX
  \usepackage[utf8]{inputenc}
  \usepackage[T2A]{fontenc}
  \usepackage{tempora}
\else
  \usepackage{fontspec}
  \IfFontExistsTF{Times New Roman}
    {\setmainfont{Times New Roman}}
    {\setmainfont{STIX Two Text}}
  \IfFontExistsTF{DejaVu Sans Mono}
    {\setmonofont{DejaVu Sans Mono}[Scale=0.82]}
    {\setmonofont{DejaVuSansMono}[Extension=.ttf, BoldFont=*-Bold,
                                  ItalicFont=*-Oblique, Scale=0.82]}
\fi
\usepackage[russian]{babel}

% Поля: слева 30 мм, справа 15 мм, сверху и снизу 20 мм
\usepackage{geometry}
\geometry{left=30mm, right=15mm, top=20mm, bottom=20mm}

\usepackage{setspace}
\singlespacing

\pretolerance=1000
\tolerance=2000
\emergencystretch=3em

\usepackage{indentfirst}
\setlength{\parindent}{1.25cm}

% Номер страницы — снизу справа
\usepackage{fancyhdr}
\pagestyle{fancy}
\fancyhf{}
\fancyfoot[R]{\thepage}
\renewcommand{\headrulewidth}{0pt}
\fancypagestyle{plain}{%
  \fancyhf{}%
  \fancyfoot[R]{\thepage}%
  \renewcommand{\headrulewidth}{0pt}%
}

\usepackage{amsmath,amssymb}
\usepackage{graphicx}
\usepackage{float}
\usepackage{booktabs}
\usepackage{tabularx}
\usepackage{multirow}
\usepackage{longtable}
\usepackage{array}

\usepackage{caption}
\captionsetup[figure]{labelsep=endash, justification=centering,
                      font={normalsize}, name=Рисунок}
\captionsetup[table]{labelsep=endash, justification=raggedright,
                     singlelinecheck=false, font={normalsize}, name=Таблица}

\usepackage{titlesec}
\titleformat{\section}
  {\newpage\normalfont\fontsize{14}{18}\selectfont\bfseries}
  {\thesection}{0.5em}{\MakeUppercase}
\titlespacing*{\section}{1.25cm}{0pt}{14pt}
\titleformat{\subsection}
  {\normalfont\fontsize{14}{18}\selectfont\bfseries}
  {\thesubsection}{0.5em}{}
\titlespacing*{\subsection}{1.25cm}{14pt}{14pt}
\titleformat{\subsubsection}
  {\normalfont\fontsize{14}{18}\selectfont\bfseries}
  {\thesubsubsection}{0.5em}{}
\titlespacing*{\subsubsection}{1.25cm}{14pt}{14pt}

\newcommand{\sectioncenter}[1]{%
  \newpage
  \phantomsection
  \addcontentsline{toc}{section}{#1}
  \begin{center}\textbf{\MakeUppercase{#1}}\end{center}
  \vspace{14pt}
}

\usepackage{tocloft}
\renewcommand{\cftsecfont}{\normalfont}
\renewcommand{\cftsecpagefont}{\normalfont}
\renewcommand{\cftsubsecfont}{\normalfont}
\renewcommand{\cftsubsecpagefont}{\normalfont}
\renewcommand{\cftsecleader}{\cftdotfill{\cftdotsep}}
\renewcommand{\cftdotsep}{1}
% babel переопределяет заголовки после преамбулы, поэтому правим через \addto
\addto\captionsrussian{\renewcommand{\contentsname}{\hfill\textbf{СОДЕРЖАНИЕ}\hfill}}
\renewcommand{\cftaftertoctitle}{\hfill}

\usepackage{enumitem}
\setlist[itemize]{noitemsep, topsep=0pt, leftmargin=1.25cm, label=--}
\setlist[enumerate]{noitemsep, topsep=0pt, leftmargin=1.25cm}

% Листинги кода и вывода программы
\usepackage{listings}
\usepackage{fancyvrb}
\lstset{
  basicstyle=\ttfamily\footnotesize,
  frame=single,
  breaklines=true,
  showstringspaces=false,
  inputencoding=utf8,
  extendedchars=true,
  columns=fullflexible,
  keepspaces=true,
  captionpos=t
}
\ifPDFTeX
\lstset{
  literate={а}{{\cyra}}1 {б}{{\cyrb}}1 {в}{{\cyrv}}1 {г}{{\cyrg}}1
           {д}{{\cyrd}}1 {е}{{\cyre}}1 {ж}{{\cyrzh}}1 {з}{{\cyrz}}1
           {и}{{\cyri}}1 {й}{{\cyrishrt}}1 {к}{{\cyrk}}1 {л}{{\cyrl}}1
           {м}{{\cyrm}}1 {н}{{\cyrn}}1 {о}{{\cyro}}1 {п}{{\cyrp}}1
           {р}{{\cyrr}}1 {с}{{\cyrs}}1 {т}{{\cyrt}}1 {у}{{\cyru}}1
           {ф}{{\cyrf}}1 {х}{{\cyrh}}1 {ц}{{\cyrc}}1 {ч}{{\cyrch}}1
           {ш}{{\cyrsh}}1 {щ}{{\cyrshch}}1 {ъ}{{\cyrhrdsn}}1
           {ы}{{\cyrery}}1 {ь}{{\cyrsftsn}}1 {э}{{\cyrerev}}1
           {ю}{{\cyryu}}1 {я}{{\cyrya}}1 {ё}{{\cyryo}}1
           {А}{{\CYRA}}1 {Б}{{\CYRB}}1 {В}{{\CYRV}}1 {Г}{{\CYRG}}1
           {Д}{{\CYRD}}1 {Е}{{\CYRE}}1 {Ж}{{\CYRZH}}1 {З}{{\CYRZ}}1
           {И}{{\CYRI}}1 {Й}{{\CYRISHRT}}1 {К}{{\CYRK}}1 {Л}{{\CYRL}}1
           {М}{{\CYRM}}1 {Н}{{\CYRN}}1 {О}{{\CYRO}}1 {П}{{\CYRP}}1
           {Р}{{\CYRR}}1 {С}{{\CYRS}}1 {Т}{{\CYRT}}1 {У}{{\CYRU}}1
           {Ф}{{\CYRF}}1 {Х}{{\CYRH}}1 {Ц}{{\CYRC}}1 {Ч}{{\CYRCH}}1
           {Ш}{{\CYRSH}}1 {Щ}{{\CYRSHCH}}1 {Ъ}{{\CYRHRDSN}}1
           {Ы}{{\CYRERY}}1 {Ь}{{\CYRSFTSN}}1 {Э}{{\CYREREV}}1
           {Ю}{{\CYRYU}}1 {Я}{{\CYRYA}}1 {Ё}{{\CYRYO}}1
           {—}{{---}}1 {–}{{--}}1 {…}{{\ldots}}1
}
\fi

% Подписи листингов — в том же стиле, что таблицы: «Листинг X.Y — Название»
\captionsetup[lstlisting]{labelsep=endash, justification=raggedright,
                          singlelinecheck=false, font={normalsize}}
\renewcommand{\lstlistingname}{Листинг}
\AtBeginDocument{\numberwithin{lstlisting}{section}}

\lstdefinestyle{python}{language=Python, keywordstyle=\bfseries}
\lstdefinestyle{shell}{language=bash, keywordstyle=\mdseries}
\lstdefinestyle{plain}{language={}, basicstyle=\ttfamily\scriptsize}

% Абзац после листинга должен сохранять отступ 1,25 см (СТП 2.1.1)
\makeatletter
\AtBeginDocument{\let\@endpetrue\@endpefalse}
\makeatother

% Подпись к блоку verbatim: \captionof вне плавающего окружения
% глобально обнуляет \parindent, поэтому отступ восстанавливается вручную
\newcommand{\listingcaption}[2]{%
  \captionof{lstlisting}{#1}\label{#2}%
  \setlength{\parindent}{1.25cm}%
}

\usepackage{hyperref}
\hypersetup{colorlinks=true, linkcolor=black, urlcolor=blue, citecolor=black}

\numberwithin{figure}{section}
\numberwithin{table}{section}
\numberwithin{equation}{section}

% Титульный лист отчёта по лабораторной работе
\newcommand{\makelabtitle}[2]{%
  \begin{titlepage}
  \thispagestyle{empty}
  \begin{center}

  МИНИСТЕРСТВО ОБРАЗОВАНИЯ РЕСПУБЛИКИ БЕЛАРУСЬ

  \vspace{0.5em}
  Учреждение образования

  \vspace{0.5em}
  \textbf{<<Белорусский государственный университет информатики и радиоэлектроники>>}

  \vspace{2em}
  Факультет компьютерных систем и сетей

  \vspace{0.5em}
  Кафедра программного обеспечения информационных технологий
  \end{center}

  \vspace{3em}
  \begin{center}
  \textbf{ОТЧЕТ}

  \vspace{0.5em}
  по лабораторной работе №#1

  \vspace{0.5em}
  по дисциплине

  \vspace{0.5em}
  <<Информационные и технические средства защиты авторских прав>>

  \vspace{0.5em}
  <<#2>>
  \end{center}

  \vspace{5em}
  \begin{flushright}
  \begin{minipage}{0.5\textwidth}
  \textbf{Выполнил:}\\
  Лебедевич Артем Владимирович\\
  магистрант группы 556301

  \vspace{0.6em}
  \textbf{Преподаватель:}\\
  Ярмолик Вячеслав Николаевич\\
  доктор технических наук\\
  профессор кафедры ПОИТ
  \end{minipage}
  \end{flushright}

  \vfill
  \begin{center}Минск, 2026~г.\end{center}
  \end{titlepage}
}

% Титульный лист реферата
\newcommand{\makerefattitle}[1]{%
  \begin{titlepage}
  \thispagestyle{empty}
  \begin{center}

  МИНИСТЕРСТВО ОБРАЗОВАНИЯ РЕСПУБЛИКИ БЕЛАРУСЬ

  \vspace{0.5em}
  Учреждение образования

  \vspace{0.5em}
  \textbf{<<Белорусский государственный университет информатики и радиоэлектроники>>}

  \vspace{2em}
  Факультет компьютерных систем и сетей

  \vspace{0.5em}
  Кафедра программного обеспечения информационных технологий
  \end{center}

  \vspace{3em}
  \begin{center}
  \textbf{РЕФЕРАТ}

  \vspace{0.5em}
  по дисциплине

  \vspace{0.5em}
  <<Информационные и технические средства защиты авторских прав>>

  \vspace{1.5em}
  на тему

  \vspace{0.5em}
  \textbf{<<#1>>}
  \end{center}

  \vspace{5em}
  \begin{flushright}
  \begin{minipage}{0.5\textwidth}
  \textbf{Выполнил:}\\
  Лебедевич Артем Владимирович\\
  магистрант группы 556301

  \vspace{0.6em}
  \textbf{Преподаватель:}\\
  Ярмолик Вячеслав Николаевич\\
  доктор технических наук\\
  профессор кафедры ПОИТ
  \end{minipage}
  \end{flushright}

  \vfill
  \begin{center}Минск, 2026~г.\end{center}
  \end{titlepage}
}


\begin{document}

\makelabtitle{3}{Кредитные операции с физическими лицами}

\tableofcontents

\section{Цель работы}

Реализовать микросервис кредитных операций с физическими лицами: справочник кредитных
программ, форму заключения кредитного договора с контролем вводимых данных, открытие счетов
клиента по плану счетов, график платежей для двух способов погашения, проводки по схеме
<<Кредитная программа>> с использованием фонда развития банка и участие в процедуре
<<Закрытие банковского дня>> вместе с депозитными программами предыдущей работы.

\section{Задание}

Вариант №\,4 --- банк <<Решение>>. Выбраны две кредитные программы для физических лиц
(таблица~\ref{tab:products}); условия приближены к действующим предложениям банка. Условный
банк в учебной системе называется BankEt, программы банка <<Решение>> служат образцом для
его справочника видов кредита.
Банки Беларуси кредитуют физических лиц в белорусских рублях, поэтому обе программы
рублёвые; справочник валют в форме договора при этом используется.

\begin{table}[H]
\caption{Кредитные программы по образцу банка <<Решение>>}
\label{tab:products}
\centering
\small
\begin{tabularx}{\textwidth}{lXX}
\toprule
 & <<R-деньги mix>> & <<R-Онлайн>> \\
\midrule
Способ погашения & ежемесячно равными аннуитетными платежами & ежемесячно проценты, вся сумма кредита в конце срока \\
Ставка & 17,65\,\% годовых & 17,65\,\% годовых \\
Сумма кредита & 500--50\,000 BYN & 300--15\,000 BYN \\
Срок договора & 6--60 месяцев & 3--36 месяцев \\
Счёт основной суммы & 2400 & 2400 \\
Счёт процентов & 2470 & 2470 \\
\bottomrule
\end{tabularx}
\end{table}

Требования задания:
\begin{enumerate}
\item внести в справочник <<План счетов>> счета, необходимые для обслуживания кредитных
линий;
\item разработать форму заключения кредитного договора (вид кредита и валюта из
справочников, номер договора, даты, срок, сумма, процент) с контролем корректности данных;
\item при заключении договора создавать не менее двух счетов --- для основной суммы и
для процентов по кредиту;
\item использовать счёт фонда развития банка из предыдущей работы со стартовым
капиталом 100~млрд белорусских рублей;
\item сформировать график платежей;
\item разработать процедуру <<Закрытие банковского дня>> и отчёт о состоянии счетов; при
двух активных кредитных программах в нём должно быть не менее десяти счетов с учётом
депозитных.
\end{enumerate}

\section{Проектирование}

\subsection{Место сервиса кредитов в системе}

Для кредитных операций добавлен четвёртый микросервис, \texttt{credit-service}
(рисунок~\ref{fig:services}). Он устроен так же, как сервис депозитов: хранит программы и
договоры в собственной базе, клиента запрашивает у модуля <<Клиенты>>, а счета и проводки
ведёт через главную книгу. Код существующих сервисов при этом не меняется: главная книга
открывает счета любого балансового кода и оповещает о новом банковском дне всех участников
из списка, поэтому подключение кредитов сводится к двум строкам справочника <<План счетов>>
и одной записи в настройке списка участников.

\begin{figure}[H]
\centering
\begin{tikzpicture}[lbl/.style={font=\scriptsize, text=black!80}]
\node[svc] (cli) at (0, 5.0) {client-service\\[-1pt]{\footnotesize клиенты (ЛР1)}\\[-3pt]{\footnotesize порт 8081}};
\node[svc] (dep) at (-5.4, 2.5) {deposit-service\\[-1pt]{\footnotesize вклады (ЛР2)}\\[-3pt]{\footnotesize порт 8083}};
\node[svc, fill=codekw!18] (crd) at (5.4, 2.5) {credit-service\\[-1pt]{\footnotesize кредиты, графики, карты}\\[-3pt]{\footnotesize порт 8084}};
\node[svc, minimum width=4.2cm] (acc) at (0, 0) {account-service\\[-1pt]{\footnotesize главная книга, банковский день}\\[-3pt]{\footnotesize порт 8082}};
\draw[flow] (dep.north) |- node[lbl, above, pos=0.75] {клиент} (cli.west);
\draw[flow] (crd.north) |- node[lbl, above, pos=0.75] {клиент} (cli.east);
\draw[flow] ([xshift=-4mm]dep.south) |- node[lbl, below, pos=0.75] {счета, проводки} ([yshift=-2.5mm]acc.west);
\draw[flow] ([xshift=4mm]crd.south) |- node[lbl, below, pos=0.75] {счета, проводки} ([yshift=-2.5mm]acc.east);
\draw[flow, dashed] ([yshift=2.5mm]acc.west) -| node[lbl, above, pos=0.25] {новый день} ([xshift=4mm]dep.south);
\draw[flow, dashed] ([yshift=2.5mm]acc.east) -| node[lbl, above, pos=0.25] {новый день} ([xshift=-4mm]crd.south);
\end{tikzpicture}
\caption{Сервис кредитов среди микросервисов банка}
\label{fig:services}
\end{figure}

\subsection{План счетов}

Для кредитных программ в план счетов добавлены два активных счёта
(таблица~\ref{tab:chart}). В отличие от депозитных, оба счёта клиента активные: их остаток
растёт по дебету и уменьшается по кредиту.

\begin{table}[H]
\caption{Счета плана счетов, используемые кредитными программами}
\label{tab:chart}
\centering
\small
\begin{tabularx}{\textwidth}{clXl}
\toprule
Код & Активность & Наименование & Назначение в работе \\
\midrule
1010 & активный & Денежные средства в кассе & касса банка \\
2400 & активный & Кредиты физическим лицам & текущий (кредитный) счёт клиента \\
2470 & активный & Начисленные процентные доходы по кредитам физическим лицам & счёт процентов по кредиту \\
7327 & пассивный & Фонд развития банка & источник кредитных средств \\
\bottomrule
\end{tabularx}
\end{table}

\subsection{Схема данных}

Схема базы сервиса кредитов приведена на рисунке~\ref{fig:er}. График платежей хранится
построчно в таблице \texttt{payment\_item}: это позволяет отмечать оплату каждого платежа и
показывать клиенту остаток долга. Таблица \texttt{card} относится к следующей работе ---
к кредитному счёту выпускается карта для банкомата.

\begin{figure}[H]
\centering
\begin{tikzpicture}
\node[entity, text width=4.6cm, anchor=north west] (product) at (0, 0) {\textbf{credit\_product}\nodepart{two} id \hfill PK\\ name, name\_en, name\_be\\ kind, currency, rate\\ min\_amount, max\_amount\\ min\_term\_months\\ max\_term\_months\\ description (ru, en, be)};
\node[entity, text width=7.0cm, anchor=north west] (contract) at (6.7, 0) {\textbf{credit\_contract}\nodepart{two} id \hfill PK\\ number \hfill UQ\\ product\_id \hfill FK\\ client\_id, client\_name \hfill $\to$ client-service\\ currency, amount, rate, term\_months\\ start\_date, end\_date, status\\ main\_account, interest\_account \hfill $\to$ account\\ principal\_paid, interest\_paid};
\node[entity, text width=4.6cm, anchor=north west] (item) at (0, -5.0) {\textbf{payment\_item}\nodepart{two} id \hfill PK\\ contract\_id, seq \hfill FK, UQ\\ due\_date\\ principal, interest\\ balance\_after\\ paid\_on};
\node[entity, text width=7.0cm, anchor=north west] (card) at (6.7, -5.2) {\textbf{card}\nodepart{two} id \hfill PK\\ number \hfill UQ\\ contract\_id \hfill FK, UQ\\ holder, pin\_hash\\ failed\_attempts, blocked};
\draw[flow] (contract.west |- product.east) -- (product.east);
\draw[flow] (item.east) -- ++(0.9, 0) |- ([yshift=-14mm]contract.west);
\draw[flow] ([xshift=10mm]card.north) -- ([xshift=10mm]card.north |- contract.south);
\end{tikzpicture}
\caption{Схема данных сервиса кредитов}
\label{fig:er}
\end{figure}

\subsection{Схема проводок}

Проводки выполняются по упрощённой схеме <<Кредитная программа>> из задания
(таблица~\ref{tab:postings}). Знак <<$+$>> означает увеличение остатка счёта, <<$-$>> ---
уменьшение: фонд развития пассивный, остальные три счёта активные.

\begin{table}[H]
\caption{Проводки процесса <<Кредитная программа>>}
\label{tab:postings}
\centering
\footnotesize
\setlength{\tabcolsep}{4.5pt}
\begin{tabular}{clcccccccc}
\toprule
 & & \multicolumn{2}{c}{Касса (А)} & \multicolumn{2}{c}{СФРБ (П)} & \multicolumn{2}{c}{Текущий счёт (А)} & \multicolumn{2}{c}{Проценты (А)} \\
\cmidrule(lr){3-4}\cmidrule(lr){5-6}\cmidrule(lr){7-8}\cmidrule(lr){9-10}
№ & Операция & Дт & Кт & Дт & Кт & Дт & Кт & Дт & Кт \\
\midrule
1 & Выделение кредита банком & & & $-$ & & $+$ & & & \\
2 & Перевод кредита в кассу & $+$ & & & & & $-$ & & \\
3 & Получение кредита через кассу & & $-$ & & & & & & \\
\midrule
4 & Внесение процентов в кассу & $+$ & & & & & & & \\
5 & Перевод процентов из кассы & & $-$ & & & & & $+$ & \\
6 & Начисление процентов банком & & & & $+$ & & & & $-$ \\
\midrule
7 & Внесение денег в кассу & $+$ & & & & & & & \\
8 & Погашение долга за кредит из кассы & & $-$ & & & $+$ & & & \\
9 & Окончание кредита & & & & $+$ & & $-$ & & \\
\bottomrule
\end{tabular}
\end{table}

Операция 1 выполняется при заключении договора. Операции 2--3 --- при получении денег
клиентом: в кассе отделения либо в банкомате (лабораторная работа №\,4), причём частями.
Остаток текущего счёта поэтому показывает, сколько кредитных средств клиент ещё может снять.
Операции 4--6 проводятся в дату уплаты процентов, 7--9 --- в дату погашения основного долга;
при аннуитетном погашении операция 9 возвращает в фонд очередную часть долга, а название
<<Окончание кредита>> носит последняя из них.

В задании операция <<Начисление процентов банком>> стоит перед внесением процентов в
кассу. В реализации проценты сначала вносятся клиентом и только затем списываются банком с
процентного счёта: итоговые обороты те же, но остаток счёта ни в какой момент не становится
отрицательным, а главная книга отрицательных остатков не допускает.

\subsection{Графики платежей}

Месячная процентная ставка равна $i = r / (12 \cdot 100)$, где $r$~--- годовая ставка.

\textbf{Аннуитет.} Ежемесячный платёж постоянен и равен
\begin{equation}
A = S \cdot \frac{i\,(1 + i)^n}{(1 + i)^n - 1},
\label{eq:annuity}
\end{equation}
где $S$~--- сумма кредита, $n$~--- срок в месяцах. В платеже номер $k$ проценты начисляются
на остаток долга $B_{k-1}$, остальное идёт в погашение:
\begin{equation}
I_k = \operatorname{round}(B_{k-1} \cdot i,\ 2), \qquad P_k = A - I_k, \qquad B_k = B_{k-1} - P_k .
\end{equation}
Последний платёж закрывает остаток целиком ($P_n = B_{n-1}$) и вбирает накопленную
погрешность округления, поэтому сумма погашений всегда в точности равна сумме кредита.

\textbf{Ежемесячная уплата процентов.} Долг не уменьшается до конца срока, поэтому каждый
месяц уплачивается одна и та же сумма $I = \operatorname{round}(S \cdot i,\ 2)$, а последний
платёж равен $S + I$.

Для кредита 12\,000 BYN на 6 месяцев под 17,65\,\% годовых $i = 0{,}0147083$,
$A = 2104{,}21$ BYN; график приведён в таблице~\ref{tab:schedule}. Тот же кредит с
ежемесячной уплатой процентов потребовал бы шести платежей по 176,50 BYN и возврата
12\,000 BYN в конце: переплата составила бы 1\,059,00 BYN против 625,27 BYN при аннуитете,
поскольку долг всё время остаётся полным.

\begin{table}[H]
\caption{Аннуитетный график платежей: 12\,000 BYN, 6 месяцев, 17,65\,\% годовых}
\label{tab:schedule}
\centering
\small
\begin{tabular}{crrrr}
\toprule
№ & Основной долг & Проценты & Платёж & Остаток долга \\
\midrule
1 & 1\,927,71 & 176,50 & 2\,104,21 & 10\,072,29 \\
2 & 1\,956,06 & 148,15 & 2\,104,21 & 8\,116,23 \\
3 & 1\,984,83 & 119,38 & 2\,104,21 & 6\,131,40 \\
4 & 2\,014,03 & 90,18 & 2\,104,21 & 4\,117,37 \\
5 & 2\,043,65 & 60,56 & 2\,104,21 & 2\,073,72 \\
6 & 2\,073,72 & 30,50 & 2\,104,22 & 0,00 \\
\midrule
Итого & 12\,000,00 & 625,27 & 12\,625,27 & \\
\bottomrule
\end{tabular}
\end{table}

\subsection{Закрытие банковского дня}

Процедура закрытия дня осталась прежней: главная книга объявляет участникам дату нового
дня и сдвигает её после успешной обработки. Участников теперь два --- сервисы депозитов и
кредитов. Сервис кредитов в ответ проводит все платежи графиков, срок которых наступил;
по условию задания клиент платит без задержки. Ключ пакета проводок включает номер платежа
(\texttt{CRD-К-000001-PAY-3}), поэтому повторная обработка дня не проведёт платёж дважды,
а после простоя сервис проведёт сразу все пропущенные платежи. Когда оплачена последняя
строка графика, договор закрывается.

\subsection{REST-интерфейс}

Операции сервиса кредитов приведены в таблице~\ref{tab:rest}.

\begin{table}[H]
\caption{REST-интерфейс сервиса кредитов}
\label{tab:rest}
\centering
\footnotesize
\begin{tabularx}{\textwidth}{llX}
\toprule
Метод & Адрес & Назначение \\
\midrule
GET & \texttt{/api/products}, \texttt{/api/meta} & виды кредитов, валюты, банковский день, номер договора \\
GET & \texttt{/api/clients} & список клиентов из модуля <<Клиенты>> \\
POST & \texttt{/api/credits/schedule} & предварительный расчёт графика платежей \\
POST & \texttt{/api/credits} & заключение договора; в ответе --- договор и PIN-конверт карты \\
GET & \texttt{/api/credits}, \texttt{/api/credits/\{id\}} & договоры; договор с графиком, счетами и проводками \\
POST & \texttt{/api/credits/\{id\}/cash} & выдача наличных с кредитного счёта через кассу \\
POST & \texttt{/api/credits/\{id\}/pin} & перевыпуск PIN-кода карты \\
POST & \texttt{/api/internal/day-close} & обработка нового банковского дня \\
POST & \texttt{/api/atm/transactions} & транзакции банкомата (лабораторная работа №\,4) \\
GET & \texttt{/api/atm/demo-cards} & демонстрационные карты (лабораторная работа №\,4) \\
\bottomrule
\end{tabularx}
\end{table}

\section{Реализация}

\subsection{Структура модуля}

Структура модуля \texttt{credit-service} показана в листинге~\ref{lst:tree}.

\listingcaption{Модуль \texttt{credit-service}}{lst:tree}
\begin{Verbatim}[fontsize=\footnotesize, frame=single, rulecolor=\color{codeframe}]
domain/        CreditProduct, CreditContract, PaymentItem, CreditKind, Card, CardNumbers
repository/    репозитории Spring Data JPA
dto/           CreditRequest, CreditView, ScheduleRow, ScheduleView, CreditIssued
service/       CreditService       заключение договора, выдача, обработка дня
               ScheduleCalculator  расчёт графиков платежей
               CreditPostings      проводки по схеме задания
               CardService         выпуск карты, PIN-код
               DemoCards           демонстрационные карты для банкомата (ЛР4)
               DayCloseProcessor   участие в закрытии банковского дня
atm/           банковская сторона протокола банкомата (ЛР4)
web/           CreditController, ReferenceController, DayCloseController, AtmController
resources/     schema.sql, data.sql (программы банка), messages*.properties, static/ (web-клиент)
\end{Verbatim}

\subsection{Расчёт графика платежей}

Расчёт вынесен в класс без состояния \texttt{ScheduleCalculator}
(листинг~\ref{lst:schedule}) и используется дважды: для предварительного показа графика в
форме и при заключении договора. Вычисления ведутся в \texttt{BigDecimal}; степень
$(1 + i)^n$ считается с точностью 16 значащих цифр, денежные суммы округляются до копеек по
правилу <<половина вверх>>.

\begin{lstlisting}[style=java, caption={Аннуитетный график (\texttt{ScheduleCalculator})}, label={lst:schedule}]
/** Аннуитетный платёж: A = S * i / (1 - (1 + i)^-n). */
public static BigDecimal annuityPayment(BigDecimal amount, BigDecimal rate, int months) {
    BigDecimal i = rate.divide(MONTHS_PERCENT, PRECISION);
    if (i.signum() == 0) {
        return amount.divide(BigDecimal.valueOf(months), 2, RoundingMode.HALF_UP);
    }
    BigDecimal growth = BigDecimal.ONE.add(i).pow(months, PRECISION);
    return amount.multiply(i).multiply(growth)
            .divide(growth.subtract(BigDecimal.ONE), 2, RoundingMode.HALF_UP);
}

private static List<ScheduleRow> annuity(BigDecimal amount, BigDecimal rate, int months, LocalDate start) {
    BigDecimal payment = annuityPayment(amount, rate, months);
    BigDecimal balance = amount;
    List<ScheduleRow> rows = new ArrayList<>();
    for (int k = 1; k <= months; k++) {
        BigDecimal interest = monthlyInterest(balance, rate);
        BigDecimal principal = k == months ? balance : payment.subtract(interest).min(balance);
        balance = balance.subtract(principal);
        rows.add(new ScheduleRow(k, start.plusMonths(k), principal, interest, balance));
    }
    return rows;
}
\end{lstlisting}

\subsection{Заключение договора}

При заключении договора сервис проверяет условия программы, открывает в главной книге
счета 2400 и 2470 на имя клиента, выделяет кредит из фонда развития, сохраняет договор с
графиком и выпускает карту (листинг~\ref{lst:open}). Если в фонде недостаточно средств,
главная книга отклоняет проводку, и её сообщение доходит до формы.

\begin{lstlisting}[style=java, caption={Счета, выделение кредита, график и карта (\texttt{CreditService.open})}, label={lst:open}]
String main = ledger.open(new OpenAccount(MAIN_CHART_CODE, currency,
        client.id().intValue(), client.fullName(), number)).number();
String interest = ledger.open(new OpenAccount(INTEREST_CHART_CODE, currency,
        client.id().intValue(), client.fullName(), number)).number();
post(key(number, "ISSUE"), number, today, List.of(CreditPostings.issue(fund(currency), main, request.amount())));

CreditContract contract = new CreditContract();
...
for (ScheduleRow row : ScheduleCalculator.build(product.getKind(), request.amount(), product.getRate(),
        request.termMonths(), today)) {
    PaymentItem item = new PaymentItem();
    ...
    contract.getSchedule().add(item);
}
contracts.saveAndFlush(contract);
PinEnvelope card = cardService.issue(contract);
return new CreditIssued(view(contract), card);
\end{lstlisting}

Проводки кредитной программы собраны в классе \texttt{CreditPostings}
(листинг~\ref{lst:postings}); платёж по графику формирует до шести операций схемы, которые
главная книга проводит одним атомарным пакетом. Как и в сервисе депозитов, название операции
передаётся кодом и переводится главной книгой при чтении на язык пользователя; сообщения
самого сервиса и названия кредитных программ также доступны на русском, английском и
белорусском языках.

\begin{lstlisting}[style=java, caption={Проводки выдачи кредита и платежа по графику (\texttt{CreditPostings})}, label={lst:postings}]
public static Posting issue(String fund, String main, BigDecimal amount) {
    // выделение кредита банком
    return Posting.of(code("posting.credit.issue"), Entry.debit(fund, amount), Entry.debit(main, amount));
}

public static List<Posting> payment(String cash, String fund, String main, String interestAccount,
                                    BigDecimal interest, BigDecimal principal, boolean last) {
    List<Posting> postings = new ArrayList<>();
    if (interest.signum() > 0) {
        // внесение процентов в кассу
        postings.add(Posting.of(code("posting.credit.interestIn"), Entry.debit(cash, interest)));
        // перевод процентов из кассы
        postings.add(Posting.of(code("posting.credit.interestFromCash"),
                Entry.credit(cash, interest), Entry.debit(interestAccount, interest)));
        // начисление процентов банком
        postings.add(Posting.of(code("posting.credit.interestCharge"),
                Entry.credit(interestAccount, interest), Entry.credit(fund, interest)));
    }
    if (principal.signum() > 0) {
        // внесение денег в кассу
        postings.add(Posting.of(code("posting.credit.cashIn"), Entry.debit(cash, principal)));
        // погашение долга за кредит из кассы
        postings.add(Posting.of(code("posting.credit.repay"), Entry.credit(cash, principal), Entry.debit(main, principal)));
        // окончание кредита — для последнего платежа, иначе возврат части кредита в фонд банка
        postings.add(Posting.of(code(last ? "posting.credit.maturity" : "posting.credit.partReturn"),
                Entry.credit(main, principal), Entry.credit(fund, principal)));
    }
    return postings;
}
\end{lstlisting}

\subsection{Контроль корректности данных договора}

Форма договора проверяется дважды: в web-клиенте до отправки и в сервисе
(таблица~\ref{tab:checks}).

\begin{table}[H]
\caption{Проверки формы кредитного договора}
\label{tab:checks}
\centering
\small
\begin{tabularx}{\textwidth}{lXc}
\toprule
Поле & Правило & Где проверяется \\
\midrule
Клиент & выбран; существует в модуле <<Клиенты>> & клиент, сервис \\
Номер договора & формат \texttt{К-000001}; не занят & клиент, сервис \\
Вид кредита & запись справочника; доступен в выбранной валюте & клиент, сервис \\
Сумма & денежная; в пределах минимума и максимума программы & клиент, сервис \\
Срок & целое число месяцев в пределах программы & клиент, сервис \\
Процент & равен ставке программы & сервис \\
Дата начала & равна текущему банковскому дню & сервис \\
Дата окончания & равна дате начала плюс срок & сервис \\
Средства банка & суммы достаточно в фонде развития & главная книга \\
\bottomrule
\end{tabularx}
\end{table}

\subsection{Обработка банковского дня}

При открытии нового дня сервис проходит по графику каждого действующего договора и
проводит неоплаченные платежи, срок которых наступил (листинг~\ref{lst:day}).

\begin{lstlisting}[style=java, caption={Платежи по графику (\texttt{CreditService.processDay})}, label={lst:day}]
for (PaymentItem item : contract.getSchedule()) {
    if (item.isPaid() || item.getDueDate().isAfter(date)) {
        continue;
    }
    boolean last = item.getSeq() == contract.getTermMonths();
    post(key(number, "PAY-" + item.getSeq()), number, date, CreditPostings.payment(cash(currency), fund(currency),
            contract.getMainAccount(), contract.getInterestAccount(), item.getInterest(), item.getPrincipal(), last));
    item.setPaidOn(date);
    contract.setInterestPaid(contract.getInterestPaid().add(item.getInterest()));
    contract.setPrincipalPaid(contract.getPrincipalPaid().add(item.getPrincipal()));
}
if (contract.getSchedule().stream().allMatch(PaymentItem::isPaid)) {
    contract.setStatus(ContractStatus.CLOSED);
    contract.setClosedOn(date);
}
\end{lstlisting}

\section{Тестирование}

Сервис кредитов покрыт 78 тестами (таблица~\ref{tab:tests}). Как и в предыдущей работе,
главная книга в тестах заменена двойником \texttt{InMemoryLedger} из тестовой части
общей библиотеки, а REST-клиенты модулей <<Клиенты>> и <<Депозиты>> --- заглушками Mockito.

\begin{table}[H]
\caption{Тесты сервиса кредитов}
\label{tab:tests}
\centering
\small
\begin{tabularx}{\textwidth}{lcX}
\toprule
Класс тестов & Тестов & Что проверяется \\
\midrule
\texttt{ScheduleCalculatorTest} & 8 & формула аннуитета, оба графика, округление \\
\texttt{CardNumbersTest} & 10 & номер карты, счётчик попыток PIN-кода \\
\texttt{CreditFlowIT} & 7 & проводки по схеме, платежи, закрытие договора \\
\texttt{CreditApiIT} & 23 & контроль данных договора, график, выдача через кассу \\
\texttt{AtmProtocolIT} & 25 & протокол банкомата (лабораторная работа №\,4) \\
\texttt{DemoCardsIT} & 5 & демонстрационные карты, ответы на трёх языках (лабораторная работа №\,4) \\
\midrule
Всего & 78 & \\
\bottomrule
\end{tabularx}
\end{table}

Тест первого платежа по аннуитетному кредиту сверяет обороты всех четырёх счетов с
таблицей~\ref{tab:postings} (листинг~\ref{lst:test}); итог прогона приведён в
листинге~\ref{lst:run}.

\begin{lstlisting}[style=java, caption={Тест платежа по графику (\texttt{CreditFlowIT})}, label={lst:test}]
@Test
@DisplayName("Аннуитет: в дату платежа проводятся проценты и часть долга, остальные дни счета не меняются")
void annuityPaymentOnDueDate() {
    CreditView credit = open("К-000001", ANNUITY, "10000.00", 6).contract();
    credits.cashOut(credit.id(), new BigDecimal("10000.00"));

    assertThat(closeDays(30)).isEmpty();                                   // 31.10 — платежей ещё нет
    List<String> events = closeDays(1);                                    // 01.11 — первый платёж
    assertThat(events).containsExactly("01.11.2026: К-000001 — платёж №1: проценты 147.08, основной долг 1606.43 BYN");

    // процентный счёт: дебет +147,08 (перевод из кассы) и кредит −147,08 (начисление банком)
    assertTurnover(credit.interestAccount(), "147.08", "147.08", "0.00");
    // текущий счёт: +10000 выдача, −10000 снятие, +1606,43 погашение из кассы, −1606,43 возврат в фонд
    assertTurnover(credit.mainAccount(), "11606.43", "11606.43", "0.00");
    // касса: транзит выданного кредита и принятого платежа 1753,51
    assertTurnover(cash, "11753.51", "11753.51", "0.00");
    // фонд: −10000 выдача, +147,08 проценты, +1606,43 возврат долга
    assertThat(ledger.balance(fund)).isEqualByComparingTo(CAPITAL.subtract(new BigDecimal("8246.49")));
}
\end{lstlisting}

\listingcaption{Результат выполнения тестов}{lst:run}
\VerbatimInput[fontsize=\fontsize{7}{8.8}\selectfont, frame=single, rulecolor=\color{codeframe}, samepage=true]{listings/tests.txt}

\section{Результаты работы}

Сценарий демонстрации: в банке действуют два вклада из предыдущей работы; два клиента
берут кредиты по разным программам; банк работает месяц. Фонд развития запущен со стартовым
капиталом 100~млрд BYN.

Контроль данных формы показан на рисунке~\ref{fig:errors}: сумма и срок выходят за
пределы программы <<R-Онлайн>>. На рисунке~\ref{fig:form} --- корректно заполненная форма с
предварительным расчётом графика: кнопка <<Рассчитать график>> запрашивает расчёт у сервиса,
не создавая договор. Числа графика совпадают с таблицей~\ref{tab:schedule}.

\screenshot{img/form_errors.png}{Контроль корректности данных кредитного договора}{fig:errors}

\screenshot[0.95\textwidth]{img/form_schedule.png}{Форма кредитного договора с графиком платежей}{fig:form}

После заключения договора открывается его карточка (рисунок~\ref{fig:open}). Кредит
зачислен на текущий счёт 2400, к счёту выпущена карта; номер карты и PIN-код показаны в
<<ПИН-конверте>> один раз --- в базе хранится только хеш PIN-кода.

\screenshot[0.95\textwidth]{img/details_open.png}{Карточка кредитного договора после заключения}{fig:open}

Первый клиент получил всю сумму через кассу, второй оставил деньги на счёте. Список
договоров приведён на рисунке~\ref{fig:list}.

\screenshot{img/list.png}{Список кредитных договоров}{fig:list}

В отчёте по счетам (рисунок~\ref{fig:report-open}) десять счетов в белорусских рублях:
касса, фонд развития, по два счёта на каждый из двух кредитных и двух депозитных договоров.
Фонд развития равен $100\,000\,000\,000 + 15\,000 - 18\,000 = 99\,999\,997\,000$ BYN: он
вырос на сумму вкладов и уменьшился на сумму выданных кредитов. На текущем счёте второго
клиента остаётся 6\,000 BYN, у первого --- ноль: дебетовый оборот 12\,000 (выдача) равен
кредитовому (снятие через кассу).

\screenshot{img/report_open.png}{Отчёт по счетам: две кредитные и две депозитные программы}{fig:report-open}

После 31 закрытия дня наступает дата первых платежей (рисунок~\ref{fig:report-month}).
По аннуитетному кредиту уплачено 176,50 BYN процентов и 1\,927,71 BYN основного долга, по
кредиту <<R-Онлайн>> --- только проценты $6000 \cdot 17{,}65 / 1200 = 88{,}25$ BYN. В том же
протоколе видны операции сервиса депозитов: оба участника обработали один и тот же день.

\screenshot{img/report_month.png}{Отчёт по счетам и протокол закрытия дня через месяц}{fig:report-month}

Карточка аннуитетного кредита после первого платежа показана на
рисунке~\ref{fig:details-month}: первая строка графика отмечена оплаченной, остаток долга
уменьшился до 10\,072,29 BYN, в проводках видны все шесть операций платежа. В карточке
кредита <<R-Онлайн>> (рисунок~\ref{fig:details-io}) основной долг не изменился.

\screenshot[0.95\textwidth]{img/details_month.png}{Аннуитетный кредит после первого платежа}{fig:details-month}

\screenshot[0.95\textwidth]{img/details_interest_only.png}{Кредит с ежемесячной уплатой процентов после первого платежа}{fig:details-io}

\sectioncenter{Заключение}

В работе реализован микросервис кредитных операций с двумя программами по образцу банка
<<Решение>>:
кредитом с аннуитетным погашением и кредитом с ежемесячной уплатой процентов и возвратом
суммы в конце срока. При заключении договора сервис проверяет условия программы, открывает
по плану счетов текущий и процентный счета клиента, выделяет кредит из фонда развития банка
и формирует график платежей; платежи проводятся в дату их наступления при закрытии
банковского дня проводками по схеме из задания.

Работа показала, что принятое ранее разделение на сервисы оправдано: новая предметная
область добавлена отдельным микросервисом без изменения кода клиентов, депозитов и главной
книги. Общие правила учёта --- расчёт сальдо, запрет отрицательного остатка, атомарность и
идемпотентность проводок --- реализованы один раз в главной книге и одинаково действуют для
депозитов и кредитов, а отчёт о состоянии счетов объединяет счета обеих программ.

Корректность расчётов и проводок подтверждена 78 автоматическими тестами.

\sectioncenter{Список использованных источников}

\begin{enumerate}
\item Об установлении Плана счетов бухгалтерского учёта в банках и небанковских
кредитно-финансовых организациях Республики Беларусь: постановление Правления Национального
банка Республики Беларусь от 29.08.2013 №\,506 [Электронный ресурс]. --- Режим доступа:
\url{https://www.nbrb.by/legislation/accountingreporting}. --- Дата доступа: 01.10.2026.
\item Кредиты на потребительские нужды~/ ЗАО <<Банк <<Решение>> [Электронный ресурс]. ---
Режим доступа: \url{https://rbank.by/life/credits/}. --- Дата доступа: 01.10.2026.
\item Четыркин, Е.~М. Финансовая математика: учебник~/ Е.~М.~Четыркин. --- 10-е изд. ---
М.: Дело, 2011. --- 392~с.
\item Ньюмен, С. Создание микросервисов~/ С.~Ньюмен. --- 2-е изд. --- СПб.: Питер, 2023. ---
624~с.
\item Spring Data JPA Reference Documentation [Электронный ресурс]. --- Режим доступа:
\url{https://docs.spring.io/spring-data/jpa/reference/}. --- Дата доступа: 01.10.2026.
\end{enumerate}

\appendixtitle{А}{Исходный код сервиса кредитов}

% Файл сгенерирован scripts/gen_listings.py — вручную не править

\subsection*{Микросервис «Кредиты» (credit-service)}

\lstinputlisting[style=xml, style=appendix, caption={credit-service/pom.xml}]{../../credit-service/pom.xml}
\lstinputlisting[style=yaml, style=appendix, caption={credit-service/src/main/resources/application.yml}]{../../credit-service/src/main/resources/application.yml}
\lstinputlisting[style=sql, style=appendix, caption={credit-service/src/main/resources/schema.sql}]{../../credit-service/src/main/resources/schema.sql}
\lstinputlisting[style=sql, style=appendix, caption={credit-service/src/main/resources/data.sql}]{../../credit-service/src/main/resources/data.sql}
\lstinputlisting[style=java, style=appendix, caption={credit-service/src/main/java/by/bsuir/bank/credit/CreditServiceApplication.java}]{../../credit-service/src/main/java/by/bsuir/bank/credit/CreditServiceApplication.java}
\lstinputlisting[style=java, style=appendix, caption={credit-service/src/main/java/by/bsuir/bank/credit/domain/Card.java}]{../../credit-service/src/main/java/by/bsuir/bank/credit/domain/Card.java}
\lstinputlisting[style=java, style=appendix, caption={credit-service/src/main/java/by/bsuir/bank/credit/domain/CardNumbers.java}]{../../credit-service/src/main/java/by/bsuir/bank/credit/domain/CardNumbers.java}
\lstinputlisting[style=java, style=appendix, caption={credit-service/src/main/java/by/bsuir/bank/credit/domain/ContractStatus.java}]{../../credit-service/src/main/java/by/bsuir/bank/credit/domain/ContractStatus.java}
\lstinputlisting[style=java, style=appendix, caption={credit-service/src/main/java/by/bsuir/bank/credit/domain/CreditContract.java}]{../../credit-service/src/main/java/by/bsuir/bank/credit/domain/CreditContract.java}
\lstinputlisting[style=java, style=appendix, caption={credit-service/src/main/java/by/bsuir/bank/credit/domain/CreditKind.java}]{../../credit-service/src/main/java/by/bsuir/bank/credit/domain/CreditKind.java}
\lstinputlisting[style=java, style=appendix, caption={credit-service/src/main/java/by/bsuir/bank/credit/domain/CreditProduct.java}]{../../credit-service/src/main/java/by/bsuir/bank/credit/domain/CreditProduct.java}
\lstinputlisting[style=java, style=appendix, caption={credit-service/src/main/java/by/bsuir/bank/credit/domain/MobileOperator.java}]{../../credit-service/src/main/java/by/bsuir/bank/credit/domain/MobileOperator.java}
\lstinputlisting[style=java, style=appendix, caption={credit-service/src/main/java/by/bsuir/bank/credit/domain/MobilePayment.java}]{../../credit-service/src/main/java/by/bsuir/bank/credit/domain/MobilePayment.java}
\lstinputlisting[style=java, style=appendix, caption={credit-service/src/main/java/by/bsuir/bank/credit/domain/PaymentItem.java}]{../../credit-service/src/main/java/by/bsuir/bank/credit/domain/PaymentItem.java}
\lstinputlisting[style=java, style=appendix, caption={credit-service/src/main/java/by/bsuir/bank/credit/repository/CardRepository.java}]{../../credit-service/src/main/java/by/bsuir/bank/credit/repository/CardRepository.java}
\lstinputlisting[style=java, style=appendix, caption={credit-service/src/main/java/by/bsuir/bank/credit/repository/CreditContractRepository.java}]{../../credit-service/src/main/java/by/bsuir/bank/credit/repository/CreditContractRepository.java}
\lstinputlisting[style=java, style=appendix, caption={credit-service/src/main/java/by/bsuir/bank/credit/repository/CreditProductRepository.java}]{../../credit-service/src/main/java/by/bsuir/bank/credit/repository/CreditProductRepository.java}
\lstinputlisting[style=java, style=appendix, caption={credit-service/src/main/java/by/bsuir/bank/credit/repository/MobileOperatorRepository.java}]{../../credit-service/src/main/java/by/bsuir/bank/credit/repository/MobileOperatorRepository.java}
\lstinputlisting[style=java, style=appendix, caption={credit-service/src/main/java/by/bsuir/bank/credit/repository/MobilePaymentRepository.java}]{../../credit-service/src/main/java/by/bsuir/bank/credit/repository/MobilePaymentRepository.java}
\lstinputlisting[style=java, style=appendix, caption={credit-service/src/main/java/by/bsuir/bank/credit/dto/CashRequest.java}]{../../credit-service/src/main/java/by/bsuir/bank/credit/dto/CashRequest.java}
\lstinputlisting[style=java, style=appendix, caption={credit-service/src/main/java/by/bsuir/bank/credit/dto/CreditDetails.java}]{../../credit-service/src/main/java/by/bsuir/bank/credit/dto/CreditDetails.java}
\lstinputlisting[style=java, style=appendix, caption={credit-service/src/main/java/by/bsuir/bank/credit/dto/CreditIssued.java}]{../../credit-service/src/main/java/by/bsuir/bank/credit/dto/CreditIssued.java}
\lstinputlisting[style=java, style=appendix, caption={credit-service/src/main/java/by/bsuir/bank/credit/dto/CreditRequest.java}]{../../credit-service/src/main/java/by/bsuir/bank/credit/dto/CreditRequest.java}
\lstinputlisting[style=java, style=appendix, caption={credit-service/src/main/java/by/bsuir/bank/credit/dto/CreditView.java}]{../../credit-service/src/main/java/by/bsuir/bank/credit/dto/CreditView.java}
\lstinputlisting[style=java, style=appendix, caption={credit-service/src/main/java/by/bsuir/bank/credit/dto/Meta.java}]{../../credit-service/src/main/java/by/bsuir/bank/credit/dto/Meta.java}
\lstinputlisting[style=java, style=appendix, caption={credit-service/src/main/java/by/bsuir/bank/credit/dto/PinEnvelope.java}]{../../credit-service/src/main/java/by/bsuir/bank/credit/dto/PinEnvelope.java}
\lstinputlisting[style=java, style=appendix, caption={credit-service/src/main/java/by/bsuir/bank/credit/dto/ScheduleRequest.java}]{../../credit-service/src/main/java/by/bsuir/bank/credit/dto/ScheduleRequest.java}
\lstinputlisting[style=java, style=appendix, caption={credit-service/src/main/java/by/bsuir/bank/credit/dto/ScheduleRow.java}]{../../credit-service/src/main/java/by/bsuir/bank/credit/dto/ScheduleRow.java}
\lstinputlisting[style=java, style=appendix, caption={credit-service/src/main/java/by/bsuir/bank/credit/dto/ScheduleView.java}]{../../credit-service/src/main/java/by/bsuir/bank/credit/dto/ScheduleView.java}
\lstinputlisting[style=java, style=appendix, caption={credit-service/src/main/java/by/bsuir/bank/credit/service/CardService.java}]{../../credit-service/src/main/java/by/bsuir/bank/credit/service/CardService.java}
\lstinputlisting[style=java, style=appendix, caption={credit-service/src/main/java/by/bsuir/bank/credit/service/CreditPostings.java}]{../../credit-service/src/main/java/by/bsuir/bank/credit/service/CreditPostings.java}
\lstinputlisting[style=java, style=appendix, caption={credit-service/src/main/java/by/bsuir/bank/credit/service/CreditService.java}]{../../credit-service/src/main/java/by/bsuir/bank/credit/service/CreditService.java}
\lstinputlisting[style=java, style=appendix, caption={credit-service/src/main/java/by/bsuir/bank/credit/service/DayCloseProcessor.java}]{../../credit-service/src/main/java/by/bsuir/bank/credit/service/DayCloseProcessor.java}
\lstinputlisting[style=java, style=appendix, caption={credit-service/src/main/java/by/bsuir/bank/credit/service/ScheduleCalculator.java}]{../../credit-service/src/main/java/by/bsuir/bank/credit/service/ScheduleCalculator.java}
\lstinputlisting[style=java, style=appendix, caption={credit-service/src/main/java/by/bsuir/bank/credit/atm/AtmConfig.java}]{../../credit-service/src/main/java/by/bsuir/bank/credit/atm/AtmConfig.java}
\lstinputlisting[style=java, style=appendix, caption={credit-service/src/main/java/by/bsuir/bank/credit/atm/AtmProcessor.java}]{../../credit-service/src/main/java/by/bsuir/bank/credit/atm/AtmProcessor.java}
\lstinputlisting[style=java, style=appendix, caption={credit-service/src/main/java/by/bsuir/bank/credit/atm/AtmRequest.java}]{../../credit-service/src/main/java/by/bsuir/bank/credit/atm/AtmRequest.java}
\lstinputlisting[style=java, style=appendix, caption={credit-service/src/main/java/by/bsuir/bank/credit/atm/AtmResponse.java}]{../../credit-service/src/main/java/by/bsuir/bank/credit/atm/AtmResponse.java}
\lstinputlisting[style=java, style=appendix, caption={credit-service/src/main/java/by/bsuir/bank/credit/atm/DepositApi.java}]{../../credit-service/src/main/java/by/bsuir/bank/credit/atm/DepositApi.java}
\lstinputlisting[style=java, style=appendix, caption={credit-service/src/main/java/by/bsuir/bank/credit/web/AtmController.java}]{../../credit-service/src/main/java/by/bsuir/bank/credit/web/AtmController.java}
\lstinputlisting[style=java, style=appendix, caption={credit-service/src/main/java/by/bsuir/bank/credit/web/CreditController.java}]{../../credit-service/src/main/java/by/bsuir/bank/credit/web/CreditController.java}
\lstinputlisting[style=java, style=appendix, caption={credit-service/src/main/java/by/bsuir/bank/credit/web/DayCloseController.java}]{../../credit-service/src/main/java/by/bsuir/bank/credit/web/DayCloseController.java}
\lstinputlisting[style=java, style=appendix, caption={credit-service/src/main/java/by/bsuir/bank/credit/web/ReferenceController.java}]{../../credit-service/src/main/java/by/bsuir/bank/credit/web/ReferenceController.java}
\lstinputlisting[style=xml, style=appendix, caption={credit-service/src/main/resources/static/index.html}]{../../credit-service/src/main/resources/static/index.html}
\lstinputlisting[style=js, style=appendix, caption={credit-service/src/main/resources/static/app.js}]{../../credit-service/src/main/resources/static/app.js}
\lstinputlisting[style=java, style=appendix, caption={credit-service/src/test/java/by/bsuir/bank/credit/AtmProtocolIT.java}]{../../credit-service/src/test/java/by/bsuir/bank/credit/AtmProtocolIT.java}
\lstinputlisting[style=java, style=appendix, caption={credit-service/src/test/java/by/bsuir/bank/credit/CardNumbersTest.java}]{../../credit-service/src/test/java/by/bsuir/bank/credit/CardNumbersTest.java}
\lstinputlisting[style=java, style=appendix, caption={credit-service/src/test/java/by/bsuir/bank/credit/CreditApiIT.java}]{../../credit-service/src/test/java/by/bsuir/bank/credit/CreditApiIT.java}
\lstinputlisting[style=java, style=appendix, caption={credit-service/src/test/java/by/bsuir/bank/credit/CreditFlowIT.java}]{../../credit-service/src/test/java/by/bsuir/bank/credit/CreditFlowIT.java}
\lstinputlisting[style=java, style=appendix, caption={credit-service/src/test/java/by/bsuir/bank/credit/ScheduleCalculatorTest.java}]{../../credit-service/src/test/java/by/bsuir/bank/credit/ScheduleCalculatorTest.java}
\lstinputlisting[style=java, style=appendix, caption={credit-service/src/test/java/by/bsuir/bank/credit/TestLedgerConfig.java}]{../../credit-service/src/test/java/by/bsuir/bank/credit/TestLedgerConfig.java}


\end{document}
